\documentclass[12pt]{article}

\usepackage{fontspec}
\usepackage{xeCJK}
\usepackage{amsmath,amssymb}
\usepackage{geometry}
\usepackage{array}
\usepackage{booktabs}

\IfFontExistsTF{Songti SC}
    {\setCJKmainfont{Songti SC}}
    {\setCJKmainfont{FandolSong-Regular}}
\IfFontExistsTF{Heiti SC}
    {\setCJKsansfont{Heiti SC}}
    {\setCJKsansfont{FandolHei-Regular}}
\IfFontExistsTF{Songti SC}
    {\setCJKmonofont{Songti SC}}
    {\setCJKmonofont{FandolSong-Regular}}

\geometry{a4paper,margin=1in}
\renewcommand{\arraystretch}{1.25}
\setlength{\parindent}{2em}

\begin{document}

\begin{center}
    {\LARGE\bfseries 编译原理作业 4\par}
    \vspace{1.0cm}
    {\large
    \begin{tabular}{rl}
    姓名：& 梁力航 \\
    学号：& 23336128 \\
    日期：& 2026年6月4日
    \end{tabular}\par}
\end{center}

\vspace{0.8cm}

\section*{第 1 题}

将表达式 $(a+b) \times c + (a+b) \times d$ 表示为三地址码的形式，
并进一步使用 common subexpression elimination (CSE) 对其进行优化。

\subsection*{原始三地址码}

\[
\begin{array}{ll}
t_1 &= a + b \\
t_2 &= t_1 \times c \\
t_3 &= a + b \\
t_4 &= t_3 \times d \\
t_5 &= t_2 + t_4
\end{array}
\]

\subsection*{CSE 优化后的三地址码}

子表达式 $a+b$ 在 $t_1$ 和 $t_3$ 中重复计算。消除公共子表达式后，$a+b$ 只计算一次，
后续直接复用 $t_1$：

\[
\begin{array}{ll}
t_1 &= a + b \\
t_2 &= t_1 \times c \\
t_3 &= t_1 \times d \\
t_4 &= t_2 + t_3
\end{array}
\]

优化效果：指令数从 5 条减少到 4 条，且减少了一次加法运算。

\section*{第 2 题}

给定数组 \texttt{int a}，按行优先的方式存放在起始于 \texttt{base} 的一片连续单元中，
数组下标从 0 开始，每个元素的字节数为 4，\texttt{a[i]} 的起始地址为：

\[
\boxed{\,base + i \times 4\,}
\]

推导：对于一维数组，按行优先（即连续存放），元素 $a[i]$ 相较于基地址偏移 $i$ 个元素。
每个元素占 4 字节，故偏移量为 $i \times 4$ 字节。

\section*{第 3 题}

以短路计算方式生成布尔表达式的中间代码时，跳转目标地址尚未确定，可使用
\boxed{\,\text{回填（backpatching）}\,}
技术避免多遍扫描。

说明：在生成短路计算的跳转指令时，某些跳转目标（如 $E.true$ 或 $E.false$）尚未生成。
回填技术通过维护一个待填地址链表，在目标地址确定后统一回填，从而只需一遍扫描即可完成中间代码生成。

\section*{第 4 题}

给定一个 $5 \times 6$ 的二维数组 $a$，$a$ 中的每个元素大小为 4 字节。
请给出 $i = a[3][2]$ 的四元式。

\subsection*{地址计算}

按行优先存储，$a[row][col]$ 的地址公式为：
\[
\text{addr}(a[row][col]) = base + (row \times n_2 + col) \times w
\]
其中 $n_2 = 6$（列数），$w = 4$（每元素字节数）。

代入 $row=3$, $col=2$：
\[
\text{addr}(a[3][2]) = base + (3 \times 6 + 2) \times 4 = base + 20 \times 4 = base + 80
\]

\subsection*{四元式序列}

四元式格式为 $(op,\ arg_1,\ arg_2,\ result)$。

\[
\begin{array}{lllll}
(1) & (*, & 3, & 6, & T_1) \\
(2) & (+, & T_1, & 2, & T_2) \\
(3) & (*, & T_2, & 4, & T_3) \\
(4) & (=[] , & a[T_3], & \_, & i)
\end{array}
\]

说明：
\begin{itemize}
    \item 第 (1) 步：计算行偏移 $row \times n_2 = 3 \times 6$，结果存入 $T_1$。
    \item 第 (2) 步：加上列号，$T_1 + 2$，结果存入 $T_2$。
    \item 第 (3) 步：乘以元素字节数 4，得到字节偏移量，存入 $T_3$。
    \item 第 (4) 步：$=[]$ 表示从数组中取值的操作，以 $a[T_3]$ 为源，结果赋给 $i$。
\end{itemize}

\end{document}
